% ============================================================
%  Chapter: When Does φ(π₁,…,πₙ)=0 Qualify as a
%           General Equation of Science?
%
%  PLACEMENT: after Chapter 8 (Form of Buckingham's Functions),
%  before Part VI (Treating With Random Variables).
%
%  REVISION NOTES (relative to the original Chapter 1 draft):
%  ─────────────────────────────────────────────────────────
%  • Introduction rewritten as a synthesis of Parts I–V,
%    not a preview of them.  Forward references to later
%    chapters are kept only where strictly necessary.
%  • A new opening section ("What we now know") grounds
%    the abstract conditions in the four concrete tools
%    the reader has already mastered.
%  • Section on the pendulum example expanded with a
%    side-by-side comparison table (model vs property-based).
%  • Section 9u.4 (Condition C6) made self-contained:
%    the additive form now references the result proved in
%    Chapter 8 rather than announcing it.
%  • Closing remarks reframed as a bridge toward the
%    probabilistic Part VI, not as an epilogue.
%  • All backward references use \ref{}; forward references
%    to copulas/Bayesian networks kept minimal.
% ============================================================
\startcontents[chapter]
\chapter{When \texorpdfstring{$\varphi(\ldots)=0$}{phi(pi)=0}
         Qualifies as a General Scientific Representation?}
\label{ch:universality_conditions}
\minitocboxed

% ─────────────────────────────────────────────────────────────────────────────
\section*{A Look Back Before Moving Forward}

At this point in the book, four powerful tools are in hand.

\begin{itemize}[leftmargin=2.5em]
  \item The \textbf{Buckingham $\pi$ theorem}
    (Chapter~\ref{ch:buckinghamPi}) guarantees that any
    physically meaningful equation among $n$ variables and $k$
    fundamental dimensions can be rewritten as a relation
    $\varphi(\pi_1,\ldots,\pi_m)=0$ among $m = n-k$
    dimensionless groups.  The theorem tells us the \emph{form}
    of the equation, but not which $\varphi$ is the right one.

  \item \textbf{Functional equations}
    (Chapter~\ref{ch:functionalequations})
    shows how to derive $\varphi$ from the structural
    \emph{properties} the phenomenon must possess---additivity,
    scale invariance, associativity, damage accumulation, and
    so on---rather than guessing or fitting it.

  \item The \textbf{ECDF transformation}
    (Chapter~\ref{ch:cdftransformation}) provides a
    non-parametric route to the same goal: by replacing every
    observation with its empirical quantile, all variables
    become dimensionless and bounded in $[0,1]$, with a
    universal interpretability scale that works across every
    domain without domain-specific knowledge.

  \item The \textbf{general functional form}
    (Chapter~\ref{ch:Piforms}) proves that, under our two natural
    aggregation conditions, the function $\varphi$ is
    \emph{forced} into the additive structure
    $\varphi(\pi_1,\ldots,\pi_n) = p\!\bigl(\sum_i p_i(\pi_i)\bigr)$,
    unifying generalized additive models, Archimedean copulas,
    and Kolmogorov--Arnold networks as special cases of a
    single axiomatically derived framework.
\end{itemize}

\noindent
Each of these tools can be used in isolation.  But together they
raise a question that is easy to overlook when one is immersed in
the technical details of any single tool: \emph{what, exactly,
does it mean for an equation to be general?}

Though we could have dealt with this question at the very beginning of the book, we have delayed it to permit the reader to have some knowledge about previous concepts that are needed to understand in depth this relevant question.

This chapter answers that question precisely.  Nine conditions are
set out---organised into three tiers of increasing depth---that
$\varphi(\pi_1,\ldots,\pi_n)=0$ must satisfy in order to deserve
the name \emph{general representation framework for scientific models}.  Readers who have
followed the development of Parts~I--V will recognise these
conditions not as new impositions but as the natural codification
of the principles that have guided every derivation in the book.

The chapter also examines how the two principal routes to
universality described above---the parametric route via functional
equations and the non-parametric route via the ECDF
transformation---satisfy the nine conditions, and where they
differ.  This comparison sets the stage for Part~VI, where the
same framework is extended to random variables and copulas.

% ─────────────────────────────────────────────────────────────────────────────
\section{The Central Question and Why It Matters}
\label{sec:ucq:question}

Throughout the history of science, equations have been proposed,
tested, abandoned, and revised.  Among this vast catalogue, a small
number have proved genuinely universal: Newton's second law,
Fourier's heat equation, Maxwell's equations, the Boltzmann entropy
relation.  What distinguishes these from the countless empirical
fits and domain-specific correlations that surround them?

One might answer: \emph{generality of application}.  But this
answer, while correct, is circular.  An equation that applies
broadly does so \emph{for a reason}---and that reason is
structural, not merely empirical.  Newton's second law applies
universally not because someone tested it in enough situations, but
because it is the necessary consequence of how force, mass, and
acceleration are \emph{conceptually related}.  The equation encodes
a logical constraint on the world, not merely a statistical
regularity.

The Buckingham $\pi$ theorem is already a major step toward
universality: it guarantees the dimensionless form.  But it does
not tell us which $\varphi$ is correct, nor whether a given
$\varphi$ is truly general or merely convenient.  Functional
equations narrow the field further: among all conceivable
functions, only those consistent with the structural properties of
the phenomenon are admissible.  The ECDF transformation opens a
third path: universal representability without committing to any
functional form at all.

The conditions set out in the following sections address the
remaining gap.  They answer the question:


%\itshape
Given a relation $\varphi(\pi_1,\ldots,\pi_n)=0$, what additional
requirements must it satisfy for us to call it a
\emph{general representation framework for scientific models} rather than a \emph{model}, a
\emph{fit}, or a \emph{convenient approximation}?

 Rather than constituting an arbitrary checklist, the nine conditions form a three-tier hierarchy that progressively transforms a dimensionless relation into a genuine scientific representation framework. Formal conditions establish mathematical legitimacy, structural conditions provide explanatory grounding, and epistemic conditions ensure interpretability and transferability. Only when all three levels are simultaneously satisfied can a relation of the form
\[
\varphi(\pi_1,\ldots,\pi_n)=0
\]
be regarded as a universal representation capable of organizing, explaining, and connecting scientific phenomena across different domains and scales.

% ─────────────────────────────────────────────────────────────────────────────
\section{Tier~1: Necessary Formal Conditions}
\label{sec:ucq:tier1}

The first tier contains conditions without which the equation
cannot even be a candidate for generality.  These are logical
prerequisites---the minimum price of admission---not marks of
quality.

\subsection{C1 --- unit invariance}

\begin{condbox}
\begin{condition}[Unit Invariance]\label{cond:ucq:unit}
The equation $\varphi(\pi_1,\ldots,\pi_n)=0$ must hold
independently of the choice of unit system: it must be invariant
under any change of the fundamental units of measurement, see \cite{FalmagneD15}.
\end{condition}
\end{condbox}

This is the foundational constraint established by Fourier and
systematised by Buckingham (Chapter~\ref{ch:buckinghamPi}): a
relation among physical quantities can be a law of nature only if
it does not depend on the arbitrary conventions humans use to
measure those quantities.  The requirement that the $\pi_i$ be
dimensionless is the precise mathematical guarantee of this
invariance.  Any equation that changes its form when one passes
from SI to Imperial units, or from degrees Celsius to Fahrenheit,
is an artifact of measurement convention, not a statement about
reality.

Condition~\ref{cond:ucq:unit} is non-negotiable.  An equation
that fails it may be useful within a fixed measurement context, but
it cannot be transferred to a new context, let alone claimed as
general.  The ECDF transformation satisfies this condition
automatically: quantiles are dimensionless order statistics,
invariant under any monotone change of units.

For example, KANs could be considered as general equations after demonstrating that they are the sets of solutions of a system of functional equations, but not because of an extremely good fit to an interpretable physical model. The reader should realize that the first condition is an step further.

\subsection{C2 --- Completeness of the $\pi$ Set}

\begin{condbox}
\begin{condition}[Completeness of the $\pi$ Set]
  \label{cond:ucq:complete}
The set $\{\pi_1,\ldots,\pi_n\}$ must capture \emph{all}
physically relevant variables of the phenomenon.  No variable
that has a non-negligible influence on the quantity of interest
may be absent from the set, and the set must be dimensionally
independent.
\end{condition}
\end{condbox}

An equation expressed in terms of an incomplete set of
dimensionless groups may appear to describe a phenomenon, but will
fail systematically outside the experimental region in which the
missing variables happen to be approximately constant.

As discussed in Chapter~\ref{ch:historicalFoundations}, this is
perhaps the most critical practical limitation of dimensional
analysis: dimensional algebra can identify all possible
dimensionless combinations of a given variable set, but it cannot
tell us whether the variable set itself is complete.  The
investigator must supply that knowledge.

Completeness also has a dual requirement: the $\pi$ set must not
include \emph{spurious} variables---quantities that have no genuine
influence on the phenomenon.  Spurious $\pi_i$ inflate the
dimensionality of the equation, obscure its structure, and can
produce numerical instabilities.  The ideal $\pi$ set is both
\emph{complete} and \emph{minimal}.

\subsection{C3 --- mathematical well-posedness}

\begin{condbox}
\begin{condition}[Mathematical Well-Posedness]
  \label{cond:ucq:wellposed}
The function $\varphi$ must be well-defined over the domain of
interest and must possess a non-trivial zero set: there must exist
points $(\pi_1,\ldots,\pi_n)$ at which $\varphi=0$ that are not
isolated singularities or artifacts of the coordinate system.
\end{condition}
\end{condbox}

This is a minimal mathematical sanity requirement.  It is rarely
violated in practice, but it is worth stating explicitly: a
relation that is identically zero, identically non-zero, or whose
zero set consists entirely of measure-zero singularities carries no
scientific content.

% ─────────────────────────────────────────────────────────────────────────────
\section{Tier~2: Structural Conditions}
\label{sec:ucq:tier2}

The second tier goes beyond formal correctness.  These conditions
distinguish a \emph{general equation} from a \emph{mere model}:
they ask not only whether the equation is formally valid, but
whether its form is \emph{justified} by something deeper than
convenience or fitting quality.

\subsection{C4 --- Derivability of the Functional Form}

\begin{condbox}
\begin{condition}[Derivability of the Functional Form]
  \label{cond:ucq:derive}
The functional form of $\varphi$ must be \emph{derivable} from known
structural properties of the phenomenon---properties such as
additivity, scale invariance, translation invariance,
associativity, order independence, or compatibility---by means of
functional or equivalent equations, rather than freely assumed or chosen on
grounds of mathematical convenience.
\end{condition}
\end{condbox}

This is, in our view, the single most important condition for
genuine universality.  It is precisely the condition that the
framework of this book has pursued from Chapter~\ref{ch:functionalequations}
onward, and it is useful to see it stated plainly. \textbf{A concrete contrast.}  Consider modelling the period $T$
of a simple pendulum with length $L$ in a gravitational field $g$. In the first approach, the power-law assumption is an educated
guess that happens to be consistent with the physics.  In the
second, it is the \emph{only} possibility consistent with the
stated properties.  These are fundamentally different epistemic
situations.  Only the second approach yields an equation that can
be called general in the deep sense: it could, in principle, be
discovered by any reasoner who correctly identifies the relevant
properties, regardless of prior knowledge of the answer.

\begin{table}[ht]
\centering
\caption{Two approaches to the pendulum period: the difference
  between a \emph{model} and a \emph{general equation}.}
\label{tab:ucq:pendulum}
\renewcommand{\arraystretch}{1.4}
\begin{tabularx}{\textwidth}{%
  >{\raggedright\arraybackslash}p{3cm}
  >{\raggedright\arraybackslash}X
  >{\raggedright\arraybackslash}X}
\toprule
\rowcolor{headerbg}
& \textbf{Model approach} & \textbf{Property-based approach} \\
\midrule
\textbf{Starting point}
  & Assume $T = C\,L^a\,g^b$ (a power law).
  & Ask: what symmetry must $T$ satisfy? \\
\textbf{How the form is obtained}
  & Exponents fixed by dimensional consistency.
  & Functional equation from dimensional invariance
    forces the form. \\
\textbf{Result}
  & $T = C\sqrt{L/g}$ --- consistent with physics,
    but the power-law was an educated guess.
  & $T \propto \sqrt{L/g}$ --- the \emph{only} form
    compatible with the stated properties. \\
\textbf{Epistemic status}
  & The form is \emph{assumed}.
    Dimensional analysis merely confirms it.
  & The form is \emph{derived}.
    Any reasoner who identifies the properties
    reaches the same conclusion. \\
\bottomrule
\end{tabularx}
\end{table}



\begin{remarkbox}
It is not sufficient that the form \emph{happens} to be derivable
after the fact.  The derivation must precede---or at least
accompany---the adoption of the form.  An equation adopted on
grounds of convenience and later shown to be derivable from
properties acquires retrospective justification, which strengthens
its claim to generality; but the derivation is what confers that
status.
\end{remarkbox}

\subsection{C5 --- Completeness Under the Given Properties}

\begin{condbox}
\begin{condition}[Completeness Under the Given Properties]
  \label{cond:ucq:unique}
The functional form of $\varphi$ must be the \emph{general}
solution (including free parameters and some arbitrary function) of the system of functional equations derived from the stated properties of the phenomenon.
\end{condition}
\end{condbox}

If many different functional forms are consistent with the stated
properties, the system of equations must contain them.

This condition has an important corollary: it makes the family of
admissible models \emph{explicit and complete}.  When the
functional equation system has a multi-parameter family of
solutions (including some arbitrary functions), we know the full set of models consistent with the
stated properties---and we know that models outside this set
violate at least one property.  This is the dual privilege
described in Chapter~\ref{ch:functionalequations}: not only do we
know what is \emph{possible}, we know what is \emph{excluded}.

In practice, uniqueness is not required. One obtains the family of solutions not a single equation.  This is not a weakness but a strength, since we keep degrees of freedom of the phenomenon---not artifacts of the fitting procedure---and the
family as a whole \emph{constitutes} the general solution, with each member corresponding to a specific instance of the phenomenon.

\subsection{C6 --- Invariance Under the Relevant Symmetry Group}

\begin{condbox}
\begin{condition}[Invariance Under the Relevant Symmetry Group]
  \label{cond:ucq:symm}
The form of $\varphi$ must be invariant under all transformations
that are physically irrelevant to the phenomenon: beyond unit
changes, these include monotone relabellings of ordinal variables,
permutations of equivalent components, and any other symmetry of
the physical system.
\end{condition}
\end{condbox}

Chapter~\ref{ch:Piforms} established one example of the most striking instance of
this condition.  Under two natural aggregation requirements---that
variables can be aggregated sequentially, and that the result must
be independent of the aggregation order---$\varphi$ is forced into
the additive form
\begin{equation}
  \varphi(\pi_1,\ldots,\pi_n)
  \;=\; p\!\Bigl(\sum_{i=1}^n p_i(\pi_i)\Bigr),
  \label{eq:ucq:additive}
\end{equation}
where $p$ is a monotone outer transformation.
This single symmetry requirement---order independence---unifies
generalized additive models, Archimedean copulas, and
Kolmogorov--Arnold networks as special cases of the same axiomatic
structure.  The result was proved in Chapter~\ref{ch:Piforms}; it
is recalled here because it is perhaps the clearest illustration
in the book of how a symmetry condition forces a functional form.

A second important instance involves ordinal variables.  The ECDF
transformation maps every variable to its quantile, and quantiles
are invariant under strictly increasing transformations of the
original scale.  This invariance property is the symmetry
exploited by Sklar's theorem and the copula framework, which will
be developed in Part~VI.

% ─────────────────────────────────────────────────────────────────────────────
\section{Tier~3: Epistemic and Interpretive Conditions}
\label{sec:ucq:tier3}

The third tier addresses questions of meaning and transferability.
An equation may be formally correct and structurally derived, yet
remain scientifically sterile if it cannot be interpreted or
transferred.  The conditions in this tier are what make a general
equation genuinely useful as a vehicle for scientific
understanding.

\subsection{C7 --- Physical Interpretability of the $\pi_i$}

\begin{condbox}
\begin{condition}[Physical Interpretability of the $\pi_i$]
  \label{cond:ucq:interp_pi}
The dimensionless groups $\pi_i$ must be interpretable in terms of
the competing mechanisms, characteristic scales, or governing
effects of the phenomenon.  They must not merely be algebraically
valid combinations of the original variables; they must carry
physical meaning that guides intuition, directs experimental
design, and enables reasoning about new situations.
\end{condition}
\end{condbox}

The classical dimensionless numbers of fluid mechanics illustrate
this requirement vividly.  The Reynolds number
$\mathrm{Re} = \rho V L/\mu$ is not merely a dimensionless ratio;
it encodes the competition between inertial and viscous forces.
The Mach number $\mathrm{Ma} = V/c$ encodes the ratio of flow
speed to acoustic propagation speed.  These numbers are
interpretable because they correspond to recognisable physical
mechanisms, and that interpretability is what allows engineers to
predict---qualitatively and quantitatively---how a system will
behave as $\mathrm{Re}$ or $\mathrm{Ma}$ change.

\begin{remarkbox}
Interpretability is a graded, not a binary, property.  Some
$\pi_i$ are immediately transparent; others become interpretable
only after the surrounding theory is developed.  It is not
required that every $\pi_i$ be interpretable at the moment the
equation is first proposed.  What is required is that
\emph{in principle} the groups admit interpretation in terms of
competing mechanisms, and that this interpretation is actively
sought and documented as the science matures.

In the ECDF setting, the interpretation is universal: the quantile
``0.02'' means ``only 2\% of observations lie below this value,''
a statement that needs no further domain knowledge.  This is a
different kind of interpretability---relative position rather than
physical mechanism---but it is real, and in many application
contexts it is exactly what is needed.
\end{remarkbox}

\subsection{C8 --- Interpretability of the Structure of $\varphi$}

\begin{condbox}
\begin{condition}[Interpretability of the Structure of $\varphi$]
  \label{cond:ucq:interp_phi}
Beyond the interpretability of the individual $\pi_i$, the
structure of $\varphi$ itself must be physically meaningful: the
equation must encode a \emph{constraint} that corresponds to
something real about the phenomenon, not merely an algebraic
identity or a curve-fitting device.
\end{condition}
\end{condbox}

Consider the S--N curve for fatigue, developed in detail in
Chapter~\ref{ch:weibull-hyperbola}.  When expressed in dimensionless
form, the percentile curves emerge as confocal rectangular
hyperbolas.  This is not a mathematical coincidence: the hyperbolic
structure reflects the existence of two physical
asymptotes---the endurance limit (stress below which no fatigue
failure occurs, regardless of the number of cycles) and the static
failure stress (stress above which failure occurs on the first
cycle).  The asymptotes are real physical phenomena; the hyperbola
is their natural geometric expression.  An investigator who
understood only the data might fit a polynomial or a spline; an
investigator who understands the physics recognises the hyperbola
as the \emph{necessary} form.

This example illustrates a general principle: when the structure
of $\varphi$ mirrors the structure of the physics, the equation
has predictive power \emph{beyond the data from which it was
inferred}.  Its extrapolations are physically grounded; its
singularities and asymptotes correspond to real phenomena; its
free parameters can be given physical interpretations.  These are
the marks of a genuinely general equation, as opposed to a
sophisticated interpolation.

\subsection{C9 --- Transferability Across Phenomena and Scales}

\begin{condbox}
\begin{condition}[Transferability Across Phenomena and Scales]
  \label{cond:ucq:transfer}
A general representation framework for scientific models must describe a \emph{class} of
phenomena, not a single instance.  The same $\varphi$ must
appear---possibly with different numerical values of its
parameters---in different physical systems that share the same
underlying structure.
\end{condition}
\end{condbox}

This is the deepest criterion: \emph{transferability} is what
separates a \emph{law} from a \emph{formula}.

The functional equations
\begin{equation}
  f(x,\,y_1+y_2) = f(x,\,y_1)+f(x,\,y_2),
  \qquad
  f(x_1+x_2,\,y)  = f(x_1,\,y)+f(x_2,\,y),
  \label{eq:ucq:bilinear}
\end{equation}
describe both the area of a rectangle and the amount of money
received under simple interest---two phenomena that appear
entirely unrelated.  The unique solution $f(x,y)=Cxy$ (Cauchy's
bilinear form) is the general equation for any relationship in
which two independently extensive quantities combine
multiplicatively.  Its transferability across domains is not
accidental; it follows from the fact that the underlying structural
property---bilinearity---is common to many phenomena.

The requirement of transferability has a practical consequence:
when proposing a general equation, one should actively seek---and
ideally demonstrate---that the same functional form appears in at
least two different, independently motivated domains.  An equation
that has been applied in one domain and claimed to be general, but
has never been transferred to any other domain, has not yet earned
the designation.




% ─────────────────────────────────────────────────────────────────────────────
\section{Two Routes to Universality: Parametric and Non-Parametric}
\label{sec:ucq:tworoutes}

Having established the nine conditions, it is instructive to
examine how the two routes studied in this book satisfy them, and
where they differ.

\subsection{The Parametric Route: Functional Equations and Derived Families}
\label{ssec:ucq:parametric}

The parametric route proceeds from \emph{properties to equations
to families}.  Starting from physically motivated
properties---additivity, max-stability, damage accumulation,
reproductive behaviour---one formulates a system of functional
equations and solves for the family of functions consistent with
those properties.  The result is a \emph{parametric family}: the
Weibull family, the generalised extreme value family, the gamma
family, the family of Archimedean copulas.

\begin{itemize}[leftmargin=2.5em]
  \item \textbf{C1~Unit invariance}: enforced by working with
    dimensionless $\pi_i$ from the outset.
  \item \textbf{C2~Completeness}: the responsibility of the physical
    analyst; the functional equation framework does not guarantee it
    but makes the gap explicit.
  \item \textbf{C3~Well-posedness}: ensured when the functional
    equation system is non-trivially solvable.
  \item \textbf{C4~Derivability}: the defining virtue of this route.
    The form is forced by the properties, not assumed.
  \item \textbf{C5~Completeness}: typically yields a parametric
    family; the family as a whole is the general equation.
  \item \textbf{C6~Symmetry}: verified by confirming that the derived
    family is invariant under the relevant transformations.
  \item \textbf{C7~Interpretability of $\pi_i$}: depends on the care
    taken in variable identification; parameters of the family
    typically carry physical interpretations.
  \item \textbf{C8~Interpretability of $\varphi$}: the structure of
    the derived family reflects the properties from which it was
    derived.
  \item \textbf{C9~Transferability}: demonstrated when the same
    family arises from the same properties in different domains.
\end{itemize}

\medskip\noindent
The strength of this route is its derivational rigour: one knows
exactly \emph{why} the functional form is what it is, see \cite{krantz1971foundations} and \cite{Narens2007}.  Its
limitation is that the free parameters and arbitrary functions of the family must be
estimated from data, and the family may not fit all instances of
the phenomenon if the underlying properties are not universally
satisfied.

\subsection{The Non-Parametric Route: The ECDF Transformation}
\label{ssec:ucq:nonparametric}

The non-parametric route achieves universality by committing to
\emph{no} functional form at all.  The ECDF transformation maps
every variable to $[0,1]$ by replacing each observation with its
empirical quantile.  The resulting variables are dimensionless,
bounded, and comparable across all domains without domain-specific
knowledge.

\begin{itemize}[leftmargin=2.5em]
  \item \textbf{C1~Unit invariance}: achieved because quantiles are
    dimensionless order statistics, invariant under any monotone
    change of units.
  \item \textbf{C6~Symmetry}: the relevant symmetry is invariance
    under monotone marginal transformations; this is what the ECDF
    transformation ensures, and it is the symmetry exploited by
    Sklar's theorem, to be developed in
    Chapter~\ref{ch:copulas}.
\end{itemize}

\medskip\noindent
Where the non-parametric route differs most significantly from the
parametric route is in Conditions~C4, C5, and C8.  The ECDF route
does not derive a specific functional form; it offers instead a
universal \emph{representation space}---the unit hypercube
$[0,1]^n$---in which any joint distribution can be studied through
its copula.  The functional form of the copula is not derived from
first principles, but must be either estimated non-parametrically
or chosen from a parametric family (at which point the parametric
route re-enters).

Condition~C7 takes a different character in the non-parametric
setting.  Individual $\pi_i$ become quantiles: ``only 2\% of
observations lie below this value.''  This is a universal
interpretability scale---applicable in medicine, engineering,
finance, and social science without modification---but it is a
scale of \emph{relative position}, not of \emph{physical
mechanism}.  Whether this constitutes genuine physical
interpretability in the sense of~C7 depends on the application:
in quality control, reliability assessment, and epidemiology,
quantile interpretability is exactly what is needed; in
fundamental physics and mechanism identification, it is one step
removed from the desired physical understanding.

\begin{remarkbox}
The two routes are not competitors but \emph{complements}.  The
parametric route provides deep structural understanding of specific
phenomena; the non-parametric route provides a universal
representation that works across all phenomena without requiring
that understanding.

A mature general representation framework for scientific models typically \emph{combines}
both: a derived parametric structure for the marginals (from
functional equations), and a copula-based description of the
dependence (from the ECDF route and Sklar's theorem), working
together within the Bayesian network framework of
Chapter~\ref{ch:bayesiannetworks}.  The chapters that follow
develop exactly this synthesis.
\end{remarkbox}


Figure \ref{fig:NineConditionsEnglish} Summarizes the nine conditions classified by the three indicated levels.

\begin{figure}[h]FF
\centering
\includegraphics[width=0.5\textwidth]{NineConditionsEnglish.png}
\caption{The nine conditions classified by levels.}
\label{fig:NineConditionsEnglish}
\end{figure}─────────────────────────────────────────────────────────────────────────────

% ─────────────────────────────────────────────────────────────────────────────
\section{Summary: A Checklist for Universality}
\label{sec:ucq:summary}

Table~\ref{tab:ucq:conditions} collects the nine conditions in one
place.

\begin{table}[ht]
\centering
\caption{The nine conditions for $\varphi(\pi_1,\ldots,\pi_n)=0$
to qualify as a general representation framework for scientific models.}
\label{tab:ucq:conditions}
\renewcommand{\arraystretch}{1.35}
\begin{tabularx}{\textwidth}{%
  >{\centering\arraybackslash}p{0.6cm}
  >{\raggedright\arraybackslash}p{4.0cm}
  >{\raggedright\arraybackslash}p{2.2cm}
  >{\raggedright\arraybackslash}X}
\toprule
\rowcolor{headerbg}
\textbf{\#} & \textbf{Condition} & \textbf{Tier}
            & \textbf{Fails if\ldots} \\
\midrule
C1 & Unit invariance
   & Formal
   & Form changes under unit transformation \\
C2 & Completeness of $\pi_i$
   & Formal
   & A governing variable is absent from the set \\
C3 & Mathematical well-posedness
   & Formal
   & $\varphi$ is trivially zero, undefined, or has no
     valid zero set \\
\midrule
C4 & Derivability of the functional form
   & Structural
   & Form is assumed, not derived from physical properties \\
C5 & Completeness under properties
   & Structural
   & Multiple incompatible forms satisfy the same stated
     properties \\
C6 & Relevant symmetry invariance
   & Structural
   & Form changes under physically irrelevant transformations \\
\midrule
C7 & Interpretability of individual $\pi_i$
   & Epistemic
   & Groups are algebraic artifacts with no physical meaning \\
C8 & Interpretability of the structure of $\varphi$
   & Epistemic
   & Equation structure has no correspondence to physical
     mechanisms \\
C9 & Transferability across phenomena
   & Epistemic
   & Form applies only to the single phenomenon for which
     it was derived \\
\bottomrule
\end{tabularx}
\end{table}

Three observations deserve emphasis.

\medskip

\textbf{The conditions are cumulative in depth.}  An equation that
satisfies only the formal conditions (C1--C3) is a \emph{valid
model}: it can be used safely, but it encodes no structural insight
and may fail outside its calibration range.  One that also
satisfies the structural conditions (C4--C6) is a \emph{derived
model}: its form is justified, and its parameters carry meaning.
One that additionally satisfies the epistemic conditions (C7--C9)
to be considered as a \emph{general representation framework for scientific models}: it encodes a logical
constraint on the world, applies to a class of phenomena, and is
interpretable at every level.  The progression is not binary but
gradual, and partial satisfaction is scientifically meaningful: it
locates the equation on a spectrum from purely empirical to
genuinely universal.

\medskip

\textbf{Condition~C4 represents an epistemological inversion.}
The usual procedure in science is to propose a functional form,
fit it to data, and assess its quality by goodness-of-fit metrics.
The programme advocated in this book inverts this order: one asks
first what structural properties the phenomenon \emph{must}
possess, derives the admissible functional forms from those
properties, and only then asks which member of the admissible
family best fits the data.  This inversion is not merely
methodological preference.  It is what separates a general
equation---which encodes a logical constraint on the world---from
a sophisticated interpolation, which merely encodes a statistical
regularity in the available data.

\medskip

\textbf{Black-box machine learning and Condition~C4.}  Neural
networks and other flexible architectures can achieve extraordinary
fitting quality, but they do so by implicitly assuming that any
sufficiently flexible function is admissible.  This is the
antithesis of the property-based approach.  The critique is not
that machine learning is unhelpful---it is immensely powerful for
prediction within the data distribution---but that it does not,
by itself, yield general equations of science.  Kolmogorov--Arnold
Networks (KANs), discussed in Chapter~\ref{ch:ArchimedeanCopulas},
represent a partial exception: their architecture is justified by a
structural theorem.  But even KANs require supplementation by the
nine conditions before their outputs can be called general
equations.

% ─────────────────────────────────────────────────────────────────────────────
\section{A Worked Example: The Fatigue S--N Equation Through All Nine
  Conditions}
\label{sec:ucq:worked}

To ground the abstract conditions in a concrete case already
encountered in the book, we trace the S--N equation for fatigue
through all nine conditions.

\subsection{Setting Up the Dimensionless Groups}

The phenomenon is the fatigue life, measured in cycles, $N$, of a material specimen
subjected to a cyclic stress ranging from a minimum stress, $\sigma_m$, and a maximum stress, $\sigma_M$, assuming that there exist a limit number of cycles, $N_0$, and an endurance limit stress, $\sigma_0$, under which failure does not occur.

The relevant dimensionless variables according to the Buckinham theorem are the three dimensionless groups:
\[
  \pi_1 = \log\left(\frac{N}{N_0}\right),
  \qquad
  \pi_2 = \frac{\sigma_M}{\sigma_0},
  \qquad
  \pi_3 = \frac{\sigma_m}{\sigma_0}.
\]

\begin{itemize}[leftmargin=2.5em]
  \item \textbf{C1~Unit invariance}: satisfied by
    construction---all three groups are dimensionless.
  \item \textbf{C2~Completeness}: for the considered test and under the assumptions stress ranging from $\sigma_m$ and $\sigma_M$ it is satisfied. If other factors as frequency, temperature, etc. other dimensionless ratios must be incorporated.
  \item \textbf{C7~Interpretability of $\pi_i$}: $\pi_1 =
    N/N_0$ is the normalised life (fraction
    of the number of cycles to failure); $\pi_2 =
    \sigma_M/\sigma_0$ and $\pi_3 =
    \sigma_m/\sigma_0$ are the normalised maximum and minimum stresses, respectively ---both are
    physically transparent stress ratios.
\end{itemize}

\subsection{Deriving the Functional Form from Physical Constraints}

To determine the form of $\varphi(\pi_1,\pi_2,\pi_3)=0$, we do
not choose a convenient curve.  Instead, we use a property leading to a functional equation that is solved and
physical constraints that must hold regardless of the material:

\begin{enumerate}[label=(\roman*)]
  \item \textbf{Limit number of cycles.} When the number of cycles $N$ tends to $N_0$ then $\sigma_M\to\infty$ i.e.\
    $\pi_2\to\infty$:
    \[
      \pi_2=\infty \quad \text{ whenever } \pi_1=\log\left(\frac{N}{N_0}\right)\to 0.
    \]
  \item \textbf{Endurance limit condition.} When the maximum applied stress tends to the endurance limit ($\sigma_M\to \sigma_0$, i.e.\
    $\pi_2\to 1$), no fatigue failure occurs:
    \[
      \pi_1\to \infty \quad \text{ whenever } \pi_2\to 1.
    \]
\end{enumerate}

These two asymptotic conditions, combined with the requirement of
strict monotonicity ($N$ decreasing in $\sigma_M$), force the
percentile curves of $N$, in the dimensionless form required by the Buckingham theorem, to be hyperplanes or of the form:
\begin{equation}
\log\left(\frac{N}{N_0}\right)\left(\frac{\sigma_M}{\sigma_0}-1\right)\frac{\sigma_m}{\sigma_0}=C(p),
  \label{eq:ucq:sn_physical}
\end{equation}
where $N_0$ and $\sigma_0$ are parameters depending on the material and the probability level $p$.

\begin{itemize}[leftmargin=2.5em]
  \item \textbf{C3~Mathematical well-posedness}: the $\varphi()$ function is defined, is not trivially zero, and has a
     valid zero set.
  \item \textbf{C4~Derivability}: the mathematical form (\ref{eq:ucq:sn_physical}) follows
    \emph{necessarily} from the two asymptotic constraints and
    monotonicity.  No assumption about the functional form was
    made; it was the only possibility.
  \item \textbf{C5~Completeness}: These solutions are the
    only compatible forms with the asymptotic properties and strict monotonicity (up to the two parameters $N_0$ and $\sigma_0$).
  \item \textbf{C6~Relevant symmetry invariance}: the equation does not depend on irrelevant (non-considered) variables.
  \item \textbf{C8~Interpretability of $\varphi$}: the horizontal
    asymptote at $\pi_2=1$ corresponds to the physical
    endurance limit; the condition $\pi_1=0$ corresponds to limit number of cycles
    failure.  The (\ref{eq:ucq:sn_physical}) structure is the geometric consequence
    of two \emph{real} physical bounds, not a modelling choice.
  \item \textbf{C9~Transferability across phenomena}: the model is applicable to other physical phenomena as demonstrated in the following section.
\end{itemize}

\subsection{Transferability: The Same Form in Three Different Domains}

The hyperbolic structure~\eqref{eq:ucq:sn_physical} appears in
physically distinct degradation mechanisms:

\begin{itemize}[leftmargin=2.5em]
  \item \textbf{Creep rupture}: time to failure as a function of
    applied stress at elevated temperature, with the creep limit
    playing the role of $\sigma_0$.
  \item \textbf{Wear and corrosion fatigue}: cumulative damage as a
    function of load cycles, where surface degradation introduces a
    shifted endurance limit.
  \item \textbf{Fracture mechanics}: the Paris--Erdogan law in
    integrated form produces S--N curves that are asymptotically
    hyperbolic in the $(\Delta K,N)$ plane.
\end{itemize}

\noindent
The same two-asymptote hyperbolic structure arising from the same
two physical constraints in three independently motivated domains
satisfies Condition~C9 and is the hallmark of a genuinely general
equation.

% ─────────────────────────────────────────────────────────────────────────────
\section{Closing Remarks: A Bridge to the Probabilistic Framework}
\label{sec:ucq:closing}

The nine conditions set out in this chapter represent a
standard---demanding but achievable---for what it means to call an
equation general.  Together they constitute a
\emph{checklist for universality}: an equation that passes all
nine has earned the right to the designation ``general representation framework for scientific models''; one that fails some of them has been identified
precisely where further work is needed.

The deepest condition is~C4: the derivability of the functional
form from structural properties of the phenomenon.  Fourier did
not \emph{choose} the heat equation because it fit his data; he
\emph{derived} it from the principle that heat flow must be
proportional to the temperature gradient.  Buckingham did not
\emph{assume} that dimensionless groups govern physical similarity;
he \emph{derived} it from the requirement that physical laws be
independent of measurement units.  The functional equations of
Chapter~\ref{ch:functionalequations} do not \emph{assume} that areas
are products of dimensions; they \emph{derive} that form from the
additivity properties of area.  In each case, the derivation is
what makes the equation a candidate to be considered as general.

\medskip

\noindent
\textbf{What comes next.}  Up to this point, the framework has
been deterministic: $\varphi(\pi_1,\ldots,\pi_n)=0$ is a
relation among definite values.  But real phenomena are
stochastic: every observation carries uncertainty, and the full
information about a system is contained not in a single value but
in a joint probability distribution over all variables.

Part~VI extends the framework to this setting.  The key insight is
that the same principles---dimensionlessness via the ECDF
transformation, functional-equation-derived structure, and order
independence---that forced the additive form in
Chapter~\ref{ch:Piforms} also determine the structure of the joint
distribution.  In particular:

\begin{itemize}[leftmargin=2.5em]
  \item The ECDF transformation produces uniform marginals on
    $[0,1]^n$, directly connecting to the definition of a copula.
  \item Sklar's theorem provides the probabilistic counterpart of
    the factorisation principle: any joint distribution decomposes
    uniquely into its marginals and a copula encoding the pure
    dependence structure.
  \item The additive form~\eqref{eq:ucq:additive} derived in
    Chapter~\ref{ch:Piforms} is precisely the form of Archimedean
    copulas, revealing that this family is not a modelling choice
    but a structural necessity under order independence.
\end{itemize}

\noindent
The general representation framework for scientific models is not discovered by searching for
the best fit among all possible functions.  It is discovered by
asking:

\begin{center}
%\itshape
What must this equation be, given what we know the phenomenon must
do?
\end{center}

\noindent
The answer to that question---deterministic or probabilistic---is
always an equation that satisfies all nine conditions.  That is,
we believe, not a coincidence.


\section{Scope and Limitations}
The results presented in this book do not imply that every scientific
law is intrinsically multilinear or can be reduced exactly to the
proposed representations.

Rather, the main claim is that a remarkably broad family of scientific,
engineering, and probabilistic models can be expressed, approximated,
or analysed within a common structural framework.

In this sense, the proposed methodology should be interpreted as a
representation language rather than as a replacement for existing
physical theories.

The role played by the proposed framework is analogous to that of
Fourier expansions. Fourier analysis does not replace the original
physical laws, yet it provides a common language for representing and
studying them. Likewise, the representations developed in this book
are intended as a unifying analytical framework rather than as a new
physical theory.

\bigskip

% ── Epigraph ──────────────────────────────────────────────────────────────────
\begin{flushright}
%\itshape
``The art of doing mathematics consists in finding\\
that special case which contains all the germs\\
of generality.''\\\medskip
\upshape --- David Hilbert
\end{flushright}
